\documentclass[12pt]{article} % default is 10pt
%%%%%%%%%%%%%%%%%%%%PREAMBLE%%%%%%%%%%%%%%%%%%%%
%%%BASIC Packages
%\usepackage[utf8x]{inputenc} % to set input encoding (not needed with XeLaTeX)
%\usepackage[T1]{fontenc} % to set font encoding in latex; difference noticed for <>
\usepackage{amsmath,amssymb,xfrac,mathrsfs,accents,bbm} % for mathematical symbols
\usepackage{mathbbol,subcaption}
%\usepackage{bbold}
%\usepackage{slashed} % for Dirac slash
%\usepackage{youngtab} % for Young tableux
%\usepackage{tipa,ulem,wasysym} % for some combining accents & other exotic symbols
%\usepackage{ulem}
\usepackage{multicol}
\usepackage{sectsty} % to change title styles
\usepackage{graphicx,enumerate, bm} % for \includegraphics command and options
%\graphicspath{{figures/}}
\usepackage{color} % to use color
\usepackage[unicode]{hyperref} % for bookmarks in Acrobat Reader and using hyperlinks
\hypersetup{bookmarksnumbered=true, bookmarksopen=true,
  bookmarksopenlevel=2, breaklinks=true, citecolor=blue,
  colorlinks=true, linkcolor=red, linktoc=page, pdfborder={0 0 0},
  pdfkeywords={Problems, }{Stochastic Calculus, }{HW},
  pdfstartview=FitH, pdfauthor={NSP}, pdfsubject={Assignment solutions
    for Mathematical Finance - Fall 2026}, pdftitle={Assignment
    solutions - Mathematical Finance}, plainpages=false, unicode=true,
  urlcolor=cyan}

%%%PAGE Styling
\usepackage{geometry} % to change the page dimensions
\geometry{letterpaper, hmargin={2.5cm,2.5cm}, vmargin={2.5cm,2cm}, headsep=0cm, headheight=1cm, footskip=1.5cm}
\usepackage{fancyhdr} % This should be set AFTER setting up the page geometry
\pagestyle{plain} % options: empty , plain , fancy
\renewcommand{\headrulewidth}{0pt}%\lhead{}\chead{}\rhead{}\lfoot{}\cfoot{\thepage}\rfoot{}
\usepackage{parskip} % needed to not screw up 'other' spacings!
\setlength{\parindent}{0.5cm}
\usepackage[titles]{tocloft} % to alter the style of the ToC
\renewcommand{\cftsecfont}{\rmfamily\mdseries\upshape}
\renewcommand{\cftsecpagefont}{\rmfamily\mdseries\upshape} % No bold!

\DeclareSymbolFontAlphabet{\ambb}{AMSb}

\newcommand{\mbb}[1]{\mathbb{#1}}
\newcommand{\mc}[1]{\mathcal{#1}}
\newcommand{\tl}[1]{\tilde{#1}}
\newcommand{\wt}[1]{\widetilde{#1}}
\newcommand{\ul}[1]{\underline{#1}}
\newcommand{\ud}{\text{d}}
\newcommand{\e}{\text{e}}
\newcommand{\mE}{\ambb{E}}
\newcommand{\mW}{\ambb{W}}
\newcommand{\mV}{\ambb{V}}
\newcommand{\mX}{\ambb{X}}
\newcommand{\mY}{\ambb{Y}}
\newcommand{\mZ}{\ambb{Z}}
\newcommand{\mR}{\ambb{R}}
\newcommand{\mTheta}{\mathbb{\Theta}}
\newcommand{\mM}{\ambb{M}}
\newcommand{\mS}{\ambb{S}}
\newcommand{\mP}{\ambb{P}}
\newcommand{\mDelta}{\mbb{\Delta}}
\newcommand{\cA}{\mathcal{A}}
\newcommand{\cF}{\mathcal{F}}

\newcommand{\Ans}{\textbf{Ans. }}


\renewcommand{\baselinestretch}{1.2}
%\numberwithin{equation}{section} % Equations numbered as <section>.#
%\renewcommand{\theequation}{\Roman{equation}}
\renewcommand{\labelenumi}{{\theenumi.}}
\renewcommand{\theenumii}{\alph{enumii}}\renewcommand{\labelenumii}{\theenumii)}
\hyphenation{hyph-ena-tion super-space anti-com-mut-a-tion}
%%%%%%%%%%%%%%%%%%%%%%%%%%%%%%%%%%%%%%%%%%%%%%%%

\begin{document}
\thispagestyle{fancy} \lhead{\bf{\large Assignment 2} }\chead{\Large
  ID5103: Mathematical Finance}\rhead{Solutions}


% \section*{On conditional expectations and sigma algebras}
\section*{}

In this assignment, we explore conditional expectations and their
computation for the random walk stochastic process. We will gradually
develop all the ingredients needed to compute conditional expectations
given the information contained in some $\sigma$-algebra. A general
sample space will be denoted $\Omega$, with more specific notation
like $\Omega_N$ reserved for sample spaces involving coin tosses. Note
that the ideas and results discussed below are guaranteed to apply
only to situations where the sample space has a finite number of
elements, and where the random variables are assumed to take a finite
number of values. The sample space and random variables of the random
walk satisfy these conditions. For other kinds of sample spaces and
random variables, one must carefully reconsider each result and check
their validity.

The random walk stochastic process $\mM_n$, $0 \leq n \leq N$, is
defined as follows. Define $\mM_0 = 0$. Then, we define the
independent random variables $\mX_i$ which depend only on the
$i^{\rm th}$ coin toss, specified by
\begin{equation}
  \mX_i(\omega_i) = \left\{\begin{array}{cc} +1 & \omega_i = H\ , \\ -1 & \omega_i = T\ . \end{array}\right.
\end{equation}
The random variable $\mM_n$ is then defined as the sum of the first
$n$ $\mX_i$:
\begin{equation}
  \mM_n(\omega_1 \cdots \omega_n) = \sum_{i=1}^n \mX_i(\omega_i)\ ,\quad\text{for}\ n = 1,\ldots,N\ .
\end{equation}

We also need the definition of an event. Simply said, an event is a
subset of the sample space. Usually, an event is a particular subset
defined by a rule or a collection of rules that the elements of the
subset must satisfy. For instance, the event $\mM_1 = 1$ is a subset
of $\Omega_N$ containing sequences of coin tosses which start with a
$H$.

Given a collection of events $A_1, A_2, \ldots, A_k$ in a sample space
$\Omega$, the $\sigma$-algebra generated by the events is the
collection of subsets obtained by taking all possible unions,
intersections and complements of $A_1, A_2, \ldots, A_k$. This
$\sigma$-algebra is written as $\sigma(A_1, \ldots, A_k)$.

Note: This $\sigma$-algebra also turns out to be the smallest
$\sigma$-algebra containing $A_1, \ldots, A_k$. That is, among all
$\sigma$-algebras containing the events $A_1, \ldots, A_k$, the
$\sigma$-algebra $\sigma(A_1, \ldots, A_k)$ is the one with the
smallest number of events in it. This notion of ``smallest
$\sigma$-algebra containing some events'' arises frequently in
probability theory.

\textbf{Exercise 1a (4 points):} Construct the $\sigma$-algebra
generated by the following subsets of $\Omega_3$:
\begin{equation}
  A_1 = \{HHH, HHT\}\ ,\quad A_2 = \{HHT, HTH, TTH\}\ .
\end{equation}
\footnote{This is an example where the subsets that generate the
  $\sigma$-algebra have non-empty intersections amongst themselves.}

\Ans We first list all $8$ elements of $\Omega_3$:
$HHH,HHT,HTH,HTT,THH,THT,TTH,TTT$. We compute the basic combinations
of $A_1$ and $A_2$:
\begin{align}
  A_1 \cap A_2 &= \{HHT\}\ , \\
  A_1^c &= \{HTH,HTT,THH,THT,TTH,TTT\}\ ,\\
  A_2^c &= \{HHH,HTT,THH,THT,TTT\}\ ,\\
  A_1 \cap A_2^c &= A_1 \setminus A_2 = \{HHH\}\ ,\\
  A_2 \cap A_1^c &= A_2 \setminus A_1 = \{HTH,TTH\}\ ,\\
  A_1^c \cap A_2^c &= (A_1 \cup A_2)^c = \{HTT,THH,THT,TTT\}\ .
\end{align}
These four sets,
\begin{equation}
  E_1 = \{HHT\}\ ,\quad E_2 = \{HHH\}\ ,\quad E_3 = \{HTH,TTH\}\ ,\quad E_4 = \{HTT,THH,THT,TTT\}\ ,
\end{equation}
are pairwise disjoint and their union is all of $\Omega_3$ (since
$1+1+2+4 = 8$), so they are exactly the atoms of $\sigma(A_1,A_2)$
(this is verified in Exercise 1b below). The full $\sigma$-algebra
$\sigma(A_1,A_2)$ consists of all $2^4 = 16$ possible unions of these
four atoms:
\begin{equation}
  \sigma(A_1,A_2) = \{\ \textstyle\bigcup_{i \in S} E_i \ :\ S \subseteq \{1,2,3,4\}\ \}\ .
\end{equation}
Explicitly, these $16$ sets are $\emptyset$, the four atoms
$E_1,E_2,E_3,E_4$, the six pairwise unions $E_1 \cup E_2 = A_1$,
$E_1 \cup E_3 = A_2$, $E_1\cup E_4$, $E_2 \cup E_3$, $E_2 \cup E_4$,
$E_3 \cup E_4$, the four triple unions $E_1\cup E_2\cup E_3 = A_1\cup
A_2$, $E_1 \cup E_2 \cup E_4$, $E_1 \cup E_3 \cup E_4 = A_1^c \cup A_2
= \Omega_3 \setminus \{HHH\}$, $E_2 \cup E_3 \cup E_4 = \Omega_3
\setminus \{HHT\}$, and finally $\Omega_3 = E_1 \cup E_2 \cup E_3 \cup
E_4$. \hfill$\blacksquare$

An atomic event or simply an atom of a $\sigma$-algebra is an event in
the $\sigma$-algebra which does not contain any other event in the
$\sigma$-algebra except itself or the empty set $\emptyset$. Suppose a
$\sigma$-algebra is generated by a finite number of subsets. For such a
$\sigma$-algebra, its atoms have the following properties:
\begin{enumerate}
\item The intersection of any two atoms is empty.
\item The union of all the atoms is the sample space itself.
\item The collection of all atoms generates the $\sigma$-algebra.
\end{enumerate}
Caution: There are $\sigma$-algebras which do not have any atomic
events in them, or even if there are atoms, they may not generate the
$\sigma$-algebra. However, we will not consider them in this course.

Note that a $\sigma$-algebra may not have been specified as being
generated by atoms but by some other collection of a finite number of
subsets. However, once we obtain the atoms of the $\sigma$-algebra, we
can show that these also generate the same $\sigma$-algebra. For
instance, let us look at the $\sigma$-algebra generated by the subsets
$C = \{HTT, HTH\}$ and $D = \{HTH, HHH\}$ of $\Omega_3$. The
$\sigma$-algebra is obtained by taking all possible unions,
intersections and complements:
\begin{gather}
  C = \{HTT, HTH\}\ ,\quad D = \{HTH, HHH\}\ ,\notag\\
  C \cup D = \{HHH, HTH, HTT\}\ ,\quad C \cap D = \{HTH\}\ ,\notag\\
  C^c = \{HHH, HHT, THH, THT, TTH, TTT\}\ ,\notag\\
  D^c = \{HHT, HTT, THH, THT, TTH, TTT\}\ ,\notag\\
  (C \cap D)^c = C^c \cup D^c = \{HHH, HHT, HTT, THH, THT, TTH, TTT\}\ ,\notag\\
  (C \cup D)^c = C^c \cap D^c = \{HHT, THH, THT, TTH, TTT\}\ ,\notag\\
  C \cup D^c = \{HHT, HTH, HTT, THH, THT, TTH, TTT\}\ ,\notag\\
  D \cup C^c = \{HHH, HHT, HTH, THH, THT, TTH, TTT\}\ ,\notag\\
  C \cap D^c = \{HTT\}\ ,\quad D \cap C^c = \{HHH\}\ ,\notag\\
  C \cup C^c = D \cup D^c = \Omega_3\ ,\quad C \cap C^c = D \cap D^c = \emptyset\ .
\end{gather}
The atoms of the above $\sigma$-algebra are
\begin{gather}
  C \cap D = \{HTH\}\ ,\quad C \cap D^c = \{HTT\}\ ,\quad D \cap C^c = \{HHH\}\, \notag\\
  C^c \cap D^c = \{HHT, THH, THT, TTH, TTT\}\ .
\end{gather}
It is easy to verify each of the three properties mentioned above for
these atoms.

\textbf{Exercise 1b (2 points):} Describe the atomic subsets of the
$\sigma$-algebra $\sigma(A_1, A_2)$ obtained in Exercise 1b.

Note that $A_1, A_2$ are not atomic subsets since they have a
non-empty intersection, and also their union is not the sample space.

\Ans Exactly as in the worked-out $C,D$ example above, we compute the
four pairwise combinations of $A_1, A_2$ and their complements. These
are precisely the four sets $E_1,E_2,E_3,E_4$ found in Exercise 1a:
\begin{gather}
  E_1 = A_1 \cap A_2 = \{HHT\}\ ,\quad E_2 = A_1 \cap A_2^c = \{HHH\}\ ,\notag\\
  E_3 = A_2 \cap A_1^c = \{HTH, TTH\}\ ,\quad E_4 = A_1^c \cap A_2^c = \{HTT, THH, THT, TTT\}\ .
\end{gather}
One checks directly that $E_1,E_2,E_3,E_4$ are pairwise disjoint, that
their union is $\Omega_3$, and that none of them contains a smaller
non-empty proper subset that itself belongs to $\sigma(A_1,A_2)$
(other than $\emptyset$). Hence $E_1,E_2,E_3,E_4$ are exactly the four
atoms of $\sigma(A_1,A_2)$. As noted in the exercise, $A_1$ and $A_2$
themselves are \emph{not} atoms: $A_1 = E_1 \cup E_2$ and
$A_2 = E_1 \cup E_3$ are each unions of two atoms, and neither equals
$\Omega_3$. \hfill$\blacksquare$

We now describe the $\sigma$-algebra generated by a random variable
$X$ which takes a finite set of values. The set of all values taken by
a random variable $X$ is called the range of $X$ and is denoted
$\text{Ran}(X)$ or $\text{Range}(X)$.

Consider subsets of sample space (events) of the form $\{X = a\}$
where $a$ is one of the possible values in $\text{Range}(X)$. For
instance, when the total number of coin tosses is $N = 4$, the event
$\{\mM_2 = 2\}$ consists of the sequences
$\{HHHH, HHHT, HHTH, HHTT\}$, i.e., all sequences starting with two
heads. Using the collection of all such subsets, form all possible
unions, intersections and complements. The resulting collection of
subsets is closed under unions, intersections and complements and is
called the $\sigma$-algebra generated by the random variable $X$ and
is denoted $\sigma(X)$.

Similarly, the $\sigma$-algebra generated by a collection of random
variables $X_1, X_2, \ldots, X_n$ is obtained as follows. Consider all
subsets of the form $\{X_1 = a_1\}, \{X_2 = a_2\}, \ldots, \{X_n =
a_n\}$ where $a_1$ runs through the range of $X_1$, $a_2$ runs through
the range of $X_2$ and so on. Form all possible unions, intersections
and complements of all these subsets: this is the $\sigma$-algebra
generated by $X_1, X_2, \ldots, X_n$ and is denoted
$\sigma(X_1, X_2, \ldots, X_n)$.

\textbf{Exercise 1c (6 points):} Compute the $\sigma$-algebras
$\sigma(\mM_1)$ and $\sigma(\mM_2)$ generated by the random variables
$\mM_1$ and $\mM_2$ respectively. You must do this separately when
$N = 3$ and when $N = 4$.

\Ans \underline{Case $N=3$:}

$\mM_1$ takes the two values $\pm 1$, with
\begin{equation}
  \{\mM_1 = 1\} = \{HHH,HHT,HTH,HTT\}\ ,\qquad \{\mM_1 = -1\} = \{THH,THT,TTH,TTT\}\ .
\end{equation}
These two sets are already complements of each other, so
\begin{equation}
  \sigma(\mM_1) = \{\emptyset,\ \{\mM_1=1\},\ \{\mM_1=-1\},\ \Omega_3\}\ ,
\end{equation}
a $\sigma$-algebra with $4$ elements.

$\mM_2$ takes the three values $2,0,-2$, with
\begin{align}
  \{\mM_2=2\} &= \{HHH,HHT\}\ ,\\
  \{\mM_2=0\} &= \{HTH,HTT,THH,THT\}\ ,\\
  \{\mM_2=-2\} &= \{TTH,TTT\}\ .
\end{align}
These three sets are already pairwise disjoint with union $\Omega_3$,
so they are the atoms of $\sigma(\mM_2)$, and
\begin{equation}
  \sigma(\mM_2) = \{\text{all } 2^3 = 8 \text{ unions of these three atoms}\}\ ,
\end{equation}
i.e.\ $\emptyset$, the three atoms above, the three sets $\{\mM_2 \neq
-2\}$, $\{\mM_2\neq 0\}$, $\{\mM_2 \neq 2\}$ obtained as pairwise
unions, and $\Omega_3$.

\underline{Case $N=4$:} Now $\Omega_4$ has $16$ elements.

$\{\mM_1=1\}$ is still just ``sequences starting with $H$'', now with
the remaining two tosses free:
\begin{equation}
  \{\mM_1=1\} = \{HHHH,HHHT,HHTH,HHTT,HTHH,HTHT,HTTH,HTTT\}\ ,
\end{equation}
an $8$-element set, and $\{\mM_1=-1\}$ is its complement (sequences
starting with $T$), also with $8$ elements. As before,
\begin{equation}
  \sigma(\mM_1) = \{\emptyset,\ \{\mM_1=1\},\ \{\mM_1=-1\},\ \Omega_4\}\ .
\end{equation}

For $\mM_2$, the value of $\mM_2$ depends only on the first two
tosses, with the last two tosses free:
\begin{align}
  \{\mM_2=2\} &= \{HHHH,HHHT,HHTH,HHTT\}\ ,\\
  \{\mM_2=0\} &= \{HTHH,HTHT,HTTH,HTTT,THHH,THHT,THTH,THTT\}\ ,\\
  \{\mM_2=-2\} &= \{TTHH,TTHT,TTTH,TTTT\}\ .
\end{align}
Again these three sets partition $\Omega_4$ and are the atoms of
$\sigma(\mM_2)$, giving a $\sigma$-algebra with $2^3=8$ elements,
exactly analogous in structure to the $N=3$ case (only the underlying
atoms are bigger, with sizes $4,8,4$ instead of $2,4,2$).
\hfill$\blacksquare$

\textbf{Exercise 1d (6 points):} Let $\cA_n$ be the $\sigma$-algebra
generated by the random variables $\mM_1$, $\mM_2, \ldots, \mM_n$. That
is, $\cA_n = \sigma(\mM_1, \mM_2, \ldots, \mM_n)$. Compute
$\cA_1 = \sigma(\mM_1)$, $\cA_2 = \sigma(\mM_1,\mM_2)$,
$\cA_3 = \sigma(\mM_1,\mM_2,\mM_3)$. For this question, you can take
the total number of coin tosses to be $N=3$.

\Ans $\cA_1 = \sigma(\mM_1)$ was already computed in Exercise 1c
above (for $N=3$):
\begin{equation}
  \cA_1 = \{\emptyset,\ \{HHH,HHT,HTH,HTT\},\ \{THH,THT,TTH,TTT\},\ \Omega_3\}\ .
\end{equation}

For $\cA_2 = \sigma(\mM_1,\mM_2)$, we need the atoms, which are the
non-empty intersections $\{\mM_1=a\}\cap\{\mM_2=b\}$:
\begin{align}
  \{\mM_1=1\}\cap\{\mM_2=2\} &= \{HHH,HHT\}\ ,\\
  \{\mM_1=1\}\cap\{\mM_2=0\} &= \{HTH,HTT\}\ ,\\
  \{\mM_1=1\}\cap\{\mM_2=-2\} &= \emptyset\ ,\\
  \{\mM_1=-1\}\cap\{\mM_2=2\} &= \emptyset\ ,\\
  \{\mM_1=-1\}\cap\{\mM_2=0\} &= \{THH,THT\}\ ,\\
  \{\mM_1=-1\}\cap\{\mM_2=-2\} &= \{TTH,TTT\}\ .
\end{align}
So there are four non-empty atoms,
$\{HHH,HHT\},\{HTH,HTT\},\{THH,THT\},\{TTH,TTT\}$, each of size $2$,
corresponding exactly to the four possible values of the pair
$(\omega_1,\omega_2)$ (namely $HH,HT,TH,TT$). This makes sense, and it must:
knowing $\mM_1$ and $\mM_2$ is exactly equivalent to knowing
$\mX_1 = \mM_1$ and $\mX_2 = \mM_2 - \mM_1$, i.e.\ to knowing the first
two coin tosses. Hence $\cA_2$ consists of all $2^4=16$ unions of
these four atoms.

For $\cA_3 = \sigma(\mM_1,\mM_2,\mM_3)$, knowing $\mM_1,\mM_2,\mM_3$ is
equivalent to knowing $\mX_1=\mM_1$, $\mX_2 = \mM_2-\mM_1$,
$\mX_3=\mM_3-\mM_2$, i.e.\ equivalent to knowing all three coin
tosses $\omega_1\omega_2\omega_3$ exactly. Hence the atoms of $\cA_3$
are the eight singletons $\{\omega\}$, $\omega \in \Omega_3$, and
$\cA_3$ is the full power set of $\Omega_3$, with $2^8 = 256$ elements.
\hfill$\blacksquare$

We next define the notion of a $\sigma$-subalgebra and that of a
filtration of $\sigma$-algebras. Given a $\sigma$-algebra $\cF$,
another $\sigma$-algebra $\mathcal G$ is said to be a
$\sigma$-subalgebra of $\cF$ if every event in $\mathcal G$ is also
contained in $\cF$. This is expressed as $\mathcal G \subseteq \cF$.

A filtration of $\sigma$-algebras is a sequence of $\sigma$-algebras
$\mathcal G_0, \mathcal G_1, \ldots, \mathcal G_n$ such that each
$\sigma$-algebra in the sequence is a $\sigma$-subalgebra of the next
one in the sequence. That is, one has the following chain of
$\sigma$-algebras $\mathcal G_0 \subseteq \mathcal G_1 \subseteq
\cdots \subseteq \mathcal G_n$.

\textbf{Exercise 1e (2 points):} Let us define $\cA_0$ to be the
trivial $\sigma$-algebra $\{\emptyset,\Omega\}$. Also recall the
$\sigma$-algebras $\cA_1, \cA_2, \cA_3$ obtained in the previous
exercise. Show that the sequence of $\sigma$-algebras
$\{\cA_n\}_{0\leq n \leq 3}$ is a filtration. That is, show that
$\cA_0 \subseteq \cA_1 \subseteq \cA_2 \subseteq \cA_3$.

Note: We frequently drop the extra qualification that $0 \leq n \leq
N$ in the subscript of the filtration and simply write $\{\cA_n\}$ for
the filtration. An $\cA_n$ without the enclosing braces will be used
to refer the $n^{\rm th}$ $\sigma$-algebra in the filtration.

\Ans $\cA_0 = \{\emptyset,\Omega_3\}$ is trivially a subset of every
$\sigma$-algebra, since every $\sigma$-algebra contains $\emptyset$ and
the whole sample space. So $\cA_0 \subseteq \cA_1$.

For $\cA_1 \subseteq \cA_2$: every element of $\cA_1$ is a union of
the two atoms $\{\mM_1=1\} = \{HHH,HHT,HTH,HTT\}$ and $\{\mM_1=-1\}
=\{THH,THT,TTH,TTT\}$. But each of these is itself a union of two of
the (finer) atoms of $\cA_2$: $\{\mM_1=1\} = \{HHH,HHT\}\cup
\{HTH,HTT\}$ and $\{\mM_1=-1\} = \{THH,THT\}\cup\{TTH,TTT\}$. Since
every element of $\cA_1$ is therefore a union of atoms of $\cA_2$, and
$\cA_2$ contains all such unions, we get $\cA_1 \subseteq \cA_2$.

For $\cA_2 \subseteq \cA_3$: each atom of $\cA_2$, e.g.
$\{HHH,HHT\}$, is itself the union of two singleton atoms of $\cA_3$,
namely $\{HHH\}\cup\{HHT\}$, and likewise for the other three atoms of
$\cA_2$. Since $\cA_3$ is the full power set of $\Omega_3$, it
certainly contains every union of atoms of $\cA_2$, so
$\cA_2 \subseteq \cA_3$.

Combining these, $\cA_0 \subseteq \cA_1 \subseteq \cA_2 \subseteq
\cA_3$, so $\{\cA_n\}_{0\leq n \leq 3}$ is indeed a filtration.
\hfill$\blacksquare$

Let us now look at conditional expectations. These are defined based
on conditional probabilities. Suppose we have events $A$ and $B$ and
we want to compute the probability that $A$ occurs given that $B$ has
occurred. This means we have to restrict our attention to that part of
the sample space in which $B$ has occurred. Recall that an event such
as $A$ or $B$ correspond to certain subsets of the sample space. For
instance, if we say $B$ is the event that the first two coin tosses
are $H$ and $T$ respectively, we have to restrict ourselves to the
subset $B = \{\omega \in \Omega_N \mid \omega_1 = H, \omega_2 = T\}$.
Similarly, an example of $A$ would be the event that the first three
coin tosses are $HTH$.

Once we restrict ourselves to the subset $B$, we then look for
occurrences of the event $A$ within the subset $B$ rather than the
entire sample space. As discussed in class, this involves the joint
occurrences of the event $A$ and $B$ which is conveyed by the
intersection $A \cap B$ which is also sometimes written $AB$. The
probability of the occurrence of $A$ given $B$ is then
\begin{equation}
  \mP\{A \mid B\} = \frac{\mP\{AB\}}{\mP\{B\}}\ .
\end{equation}
Let us consider the situation where the number of coin tosses in a
sequence is $N=3$ so that we are dealing with the sample space of
sequences of $3$ coin tosses. When $A$ is the event that the three
coin tosses are $HTH$ and $B$ is the event that the first two coin
tosses are $HT$, we have $A = \{HTH\}$, $B = \{HTH, HTT\}$, and $AB =
\{HTH\}$. We have $\mP\{B\} = p^2 q + pq^2 = pq$ and
$\mP\{AB\} = p^2 q$. Thus, we have
\begin{equation}
  \mP\{A\mid B\} = \frac{\mP\{AB\}}{\mP\{B\}} = \frac{p^2 q}{pq} = p\ .
\end{equation}
On the contrary, if $A$ is the event that the first three coin tosses
are $TTH$ and $B$ is the event that the first two coin tosses are
$HT$, these events cannot jointly occur. So, $\mP\{AB\}=0$ in this
case since $AB = \emptyset$. This gives $\mP\{A\mid B\} = 0$.

In other words, one has to recompute the probability of occurrence of
an event $A$ given new information that the event $B$ has occurred.
This new probability is the conditional probability $\mP\{A\mid B\}$
and is typically different from the original probability of
occurrence of $A$ which is simply $\mP\{A\}$.

Another example which will be useful in defining conditional
expectations is when we take the event $A$ to be a singleton set,
i.e., a set that contains a single element of the sample space
$A=\{\omega\}$. The original probability of a sequence with $n_H$
heads and $n_T$ tails is $p^{n_H}q^{n_T}$. Let us again take $B$ to be
some event in the sample space (e.g., the subset of sequences starting
with $HT$). If the sequence $\omega$ is in $B$ (that is, it starts
with $HT$), then $AB = \{\omega\}$. If $\omega$ is not in $B$, then $AB
= \emptyset$. Thus, the joint probability of $A=\{\omega\}$ and $B$
occurring is
\begin{equation}
  \mP\{\{\omega\}\cap B\} = \left\{\begin{array}{cc} \mP\{\omega\} & \omega \in B\ ,\\ 0 & \omega \notin B\ .\end{array}\right.
\end{equation}
This then gives the conditional probability
\begin{equation}
  \mP\{\omega \mid B\} = \left\{\begin{array}{cc} \dfrac{\mP\{\omega\}}{\mP\{B\}} & \omega \in B\ ,\\[2mm] 0 & \omega \notin B\ . \end{array}\right.
\end{equation}
For example, in the case of $B$ being the event where the first two
coin tosses are $HT$, we have shown earlier that $\mP\{B\} = pq$. We
also know $\mP\{\omega\} = p^{n_H(\omega)}q^{n_T(\omega)}$ where
$n_H(\omega)$ and $n_T(\omega)$ are the number of heads and tails in
the sequence $\omega$. This gives
\begin{equation}
  \mP\{\omega\mid B\} = \left\{\begin{array}{cc} p^{n_H(\omega)-1}q^{n_T(\omega)-1} & \omega \in B\ ,\\ 0 & \omega \notin B\ . \end{array}\right.
\end{equation}
Again, we see that the probability of occurrence of a sequence in
sample space is revised if we incorporate the information that the
first two coin tosses are $HT$.

The expectation value of a random variable $Y$ given the knowledge
that an event $B \subseteq \Omega_N$ has occurred is defined as
\begin{equation}
  \mE[Y\mid B] = \sum_{\omega \in \Omega_N} Y(\omega)\,\mP\{\omega\mid B\} = \sum_{\omega \in B} Y(\omega)\, \mP\{\omega\mid B\}\ .
\end{equation}
The main change compared to the formula for an ordinary expectation
value is that the probability associated to each such element $\omega$
is $\mP\{\omega \mid B\}$ and is typically different from the original
probability $\mP\{\omega\}$. Note that the summation over all elements
of the sample space can now be restricted to the subset $B$ since the
conditional probabilities for $\omega \notin B$ are zero and hence do
not contribute to the sum.

It is clear from the above formula that the conditional expectation
value is fixed once the event $B$ is specified. If the event $B$ is
changed to another event, then the expectation value may also change.

Thus, the conditional expectation value $\mE[Y\mid B]$ is a function
of the event $B$. Just like conditional probability is a revised
probability of the occurrence of an event, the conditional expectation
$\mE[Y \mid B]$ is a revised expectation value of the random variable
$Y$ given that the event $B$ has occurred.

\textbf{Exercise 1f (2 points):} Let us go back to the random walk.
Let the event $B$ be the subset of the sample space $\Omega_N$
consisting of sequences $\omega$ with $\mM_3 = 1$. That is,
\begin{equation}
  B = \{\omega \in \Omega_N \mid \mM_3(\omega) = 1\}\ ,
\end{equation}
also written in short as $B=\{\mM_3=1\}$. Compute the conditional
expectation value
\begin{equation}
  \mE[\mM_4 \mid B]\ ,\qquad\text{also written as}\qquad \mE[\mM_4\mid \mM_3=1]\ ,
\end{equation}
using the formula (11).\footnote{formula 11 in the question sheet and 38 in this solution sheet.} You should compute all the required
probabilities that appear in the formula (11). For the purposes of
this question, you can take the number of tosses to be $N=4$ and
borrow any required results from previous exercises.

\Ans With $N=4$, $B=\{\mM_3=1\}$ depends only on the first three
tosses: $\mM_3=1$ means $2$ heads and $1$ tail occurred among the first
three tosses, which happens for exactly $\binom{3}{2}=3$ arrangements,
$HHT,HTH,THH$, each with probability $p^2q$. The fourth toss is free.
So, as a subset of $\Omega_4$, $B$ consists of $6$ sequences:
\begin{align}
  HHTH\ (\mM_4=2,\ \text{prob } p^3q)\ ,&\qquad HHTT\ (\mM_4=0,\ \text{prob } p^2q^2)\ ,\\
  HTHH\ (\mM_4=2,\ \text{prob } p^3q)\ ,&\qquad HTHT\ (\mM_4=0,\ \text{prob } p^2q^2)\ ,\\
  THHH\ (\mM_4=2,\ \text{prob } p^3q)\ ,&\qquad THHT\ (\mM_4=0,\ \text{prob } p^2q^2)\ .
\end{align}
So we have $3$ sequences with $\mM_4=2$ (each of probability $p^3q$)
and $3$ sequences with $\mM_4=0$ (each of probability $p^2q^2$).

First, $\mP\{B\} = 3p^3q + 3p^2q^2 = 3p^2q(p+q) = 3p^2q$. Using
(9)-(10)\footnote{again, 9-10 in the question sheet, and 36-37 in this sheet.}, for $\omega \in B$ with $\mM_4(\omega)=2$ we get
$\mP\{\omega\mid B\} = p^3q/(3p^2q) = p/3$, and for $\omega\in B$ with
$\mM_4(\omega)=0$ we get $\mP\{\omega \mid B\} = p^2q^2/(3p^2q) = q/3$.
Then, using (11),
\begin{equation}
  \mE[\mM_4\mid B] = 3\cdot 2 \cdot \frac{p}{3} + 3 \cdot 0 \cdot \frac{q}{3} = 2p\ .
\end{equation}
(As a check: since $\mX_4$ is independent of the first three tosses,
$\mE[\mM_4\mid B] = \mE[\mM_3\mid B] + \mE[\mX_4] = 1 + (p-q) = 2p$, which agrees.) For the symmetric random walk
$p=q=\tfrac12$ this gives $\mE[\mM_4\mid B] = 1$. \hfill$\blacksquare$

As discussed above, the conditional expectation $\mE[Y\mid B]$ is a
function of the event $B$. However, we know that information present
in specifying one event $B$ in sample space can be expressed in terms
of the $\sigma$-algebra generated by the event $B$ via complements,
unions and intersections. This is in fact the $\sigma$-algebra
\begin{equation}
  \sigma(B) = \{B, B^c, \emptyset, \Omega\}\ ,
\end{equation}
and is also the smallest $\sigma$-algebra of any possible
$\sigma$-algebra that contains the event $B$. In other words, any other
$\sigma$-algebra that contains $B$ must contain more subsets than
$\sigma(B)$.

The reason to consider this extra layer of abstraction is that, if one
knows the probability $\mP\{B\}$ of the event $B$ occurring, one also
immediately knows the probability $\mP\{B^c\}$ of the event $B^c$
occurring ($\mP\{B^c\} = 1-\mP\{B\}$). Further, we know trivially that
$B \cap B^c = \emptyset$ is the null event with $\mP\{\emptyset\} =
0$, i.e., no event occurs, and $B \cup B^c = \Omega$ with $\mP\{\Omega\}
= 1$, i.e., some event occurs. Together, this $\sigma$-algebra
completely summarizes the new information that can be specified once
the event $B$ is specified, along with the probabilities -- either we
say $B$ occurs, or $B^c$ occurs, or $B\cap B^c = \emptyset$ occurs,
which is the same as the information that no sequence in the sample
space has occurred, or finally $B \cup B^c = \Omega$ occurs, which is
the same as having no prior information about any specific event
occurring.

The expectation value of some random variable $Y$ can now be computed
for every event in $\sigma(B)$ and these expectation values give us a
better understanding as to how the expected value of $Y$ is revised
given the information in $\sigma(B)$.

\textbf{Exercise 1g (4 points):} One can now compute the conditional
expectation value of a random variable $Y$ given each of the subsets in
the $\sigma$-algebra $\sigma(B)$. Take $Y=\mM_4$, $B=\{\mM_3=1\}$ as in
the previous exercise and compute
\begin{equation}
  \mE[\mM_4\mid B^c]\ ,\quad \mE[\mM_4\mid \emptyset]\ ,\quad \mE[\mM_4\mid \Omega]\ .
\end{equation}
For this question, you can take $N=4$.

\Ans \underline{$\mE[\mM_4\mid \Omega]$:} Since $\mP\{\omega\mid
\Omega\} = \mP\{\omega\}/\mP\{\Omega\} = \mP\{\omega\}$ for every
$\omega$, conditioning on the whole sample space changes nothing, so
$\mE[\mM_4\mid\Omega]$ is simply the ordinary (unconditional)
expectation of $\mM_4$:
\begin{equation}
  \mE[\mM_4\mid\Omega] = \mE[\mM_4] = 4(p-q)\ ,
\end{equation}
using the fact that $\mE[\mM_n] = n(p-q)$ for the random walk (each
step contributes $\mE[\mX_i] = p - q$ on average, independently).

\underline{$\mE[\mM_4\mid \emptyset]$:} Formula (11) defines this as a
sum over $\omega \in \emptyset$, which is an empty sum. By the usual
convention that an empty sum equals $0$,
\begin{equation}
  \mE[\mM_4\mid \emptyset] = 0\ .
\end{equation}
(Strictly speaking $\mP\{\emptyset\}=0$ makes the ratio
$\mP\{\omega\mid\emptyset\}$ in (9) an undefined $0/0$; the value $0$
above is simply the convention used to avoid quirks that arise otherwise.)

\underline{$\mE[\mM_4\mid B^c]$:} We use the law of total expectation:
since $B$ and $B^c$ partition $\Omega_4$,
\begin{equation}
  \mE[\mM_4] = \mE[\mM_4\mid B]\,\mP\{B\} + \mE[\mM_4\mid B^c]\,\mP\{B^c\}\ .
\end{equation}
{From} Exercise 1f, $\mE[\mM_4\mid B] = 2p$ and $\mP\{B\}=3p^2q$, so
$\mP\{B^c\} = 1-3p^2q$, and hence
\begin{equation}
  \mE[\mM_4\mid B^c] = \frac{\mE[\mM_4] - 2p\cdot 3p^2q}{1-3p^2q} = \frac{4(p-q) - 6p^3q}{1-3p^2q}\ .
\end{equation}
As a check, one can also compute this directly from (11) by summing
over the $10$ sequences of $B^c$: these are $HHHH\ (p^4,\mM_4{=}4)$,
$HHHT\ (p^3q,\mM_4{=}2)$, three sequences $HTTH,THTH,TTHH$ each
$(p^2q^2,\mM_4{=}0)$, four sequences $HTTT,THTT,TTHT,TTTH$ each
$(pq^3,\mM_4{=}-2)$, and $TTTT\ (q^4,\mM_4{=}-4)$; summing
$\mM_4(\omega)\mP\{\omega\}$ over these and dividing by
$\mP\{B^c\}=1-3p^2q$ reproduces the same expression above. For the
symmetric random walk $p=q=\tfrac12$, this gives
$\mE[\mM_4\mid B^c] = -\tfrac{3}{5}$. \hfill$\blacksquare$

Next, we define the conditional expectation random variable
$\mE[X\mid\cF]$ which involves a random variable $X$ and information
given in terms of a $\sigma$-algebra $\cF$. The definition we give here
holds only for a $\sigma$-algebra which is generated by its atoms. (As
mentioned earlier, not every $\sigma$-algebra has atoms, or even if it
has atoms, may not be generated by them.)

Since the atoms of the $\sigma$-algebra $\cF$ are all disjoint, it can
be shown that any given element $\omega$ of sample space belongs to
one and only one atom. First, since the union of all atoms is the
sample space, every element of sample space must belong to at least
one atom. If a given element $\omega$ belongs to more than one atom,
then the two atoms have a common element $\omega$ which means that
they are no longer disjoint. Hence, every element $\omega$ of sample
space belongs to one and only atom of the $\sigma$-algebra $\cF$. Note
that this is true for any $\sigma$-algebra generated by its atoms.

This fact allows us to define the conditional expectation random
variable $\mE[X\mid\cF]$ by giving its value on each element of sample
space. Let us call the unique atom containing a given element $\omega$
as $A^\omega$. Then, the value of $\mE[X\mid \cF]$ on $\omega$ is given
by the conditional expectation value $\mE[X\mid A^\omega]$:
\begin{equation}
  \mE[X\mid\cF](\omega) = \mE[X\mid A^\omega]\ .
\end{equation}
Note that the notation on the left and right hand sides are very
similar. On the left hand side, the conditioning is on a full
$\sigma$-algebra and this conditional expectation is always
interpreted as a random variable, i.e., a function on sample space. On
the right hand side, the conditioning is on a particular atomic event
$A^\omega$ and this is always a conditional expectation value, i.e., a
number.

For instance, suppose $B$ is an event, e.g., the event consisting of
all coin toss sequences starting with $HT$. The $\sigma$-algebra
generated by $B$ is $\{B, B^c, \emptyset, \Omega\}$. Clearly, the
events $B$ and $B^c$ are the atoms of $\sigma(B)$. Thus, the value of
the random variable $\mE[Y\mid \sigma(B)]$ on the elements
$\omega \in \Omega$ are given by
\begin{equation}
  \mE[Y\mid\sigma(B)](\omega) = \left\{\begin{array}{cc} \mE[Y\mid B] & \omega \in B\ ,\\ \mE[Y \mid B^c] & \omega \in B^c\ . \end{array}\right.
\end{equation}

\textbf{Exercise 1h (4 points):} Compute the random variables
$\mE[\mM_3\mid\sigma(\mM_2)]$ and $\mE[\mM_3\mid \cA_2]$ where
$\cA_2 = \sigma(\mM_1,\mM_2)$. How are these answers different? For
this exercise, you can take $N=3$ and use the $\sigma$-algebras for
this case.

\Ans \underline{$\mE[\mM_3\mid \sigma(\mM_2)]$:} From Exercise 1c, the
atoms of $\sigma(\mM_2)$ (for $N=3$) are $\{\mM_2=2\}=\{HHH,HHT\}$,
$\{\mM_2=0\}=\{HTH,HTT,THH,THT\}$ and $\{\mM_2=-2\}=\{TTH,TTT\}$. On
each atom, $\mM_3 = \mM_2 + \mX_3$, and since $\mX_3$ is independent of
$(\omega_1,\omega_2)$ (hence independent of the atom), we get
$\mE[\mM_3 \mid \mM_2=k] = k + \mE[\mX_3] = k+(p-q)$ for each atom
value $k$. (One can check this directly, e.g.\ for $k=0$:
$\mP\{\mM_2=0\}=2pq$ and $\sum_{\omega} \mM_3(\omega)\mP(\omega) =
p^2q - pq^2+p^2q-pq^2 = 2pq(p-q)$, giving $\mE[\mM_3\mid\mM_2=0] =
p-q$, matching $0+(p-q)$.) Hence
\begin{equation}
  \mE[\mM_3\mid\sigma(\mM_2)](\omega) = \left\{\begin{array}{cl} 2+(p-q) & \omega \in \{HHH,HHT\}\ ,\\ (p-q) & \omega \in \{HTH,HTT,THH,THT\}\ ,\\ -2+(p-q) & \omega \in \{TTH,TTT\}\ . \end{array}\right.
\end{equation}
In other words, as a random variable, $\mE[\mM_3\mid\sigma(\mM_2)] =
\mM_2 + (p-q)$.

\underline{$\mE[\mM_3\mid \cA_2]$:} From Exercise 1d, the (finer)
atoms of $\cA_2$ are $\{HHH,HHT\}$, $\{HTH,HTT\}$, $\{THH,THT\}$,
$\{TTH,TTT\}$, corresponding to $\omega_1\omega_2 = HH,HT,TH,TT$
respectively. On the atom $\omega_1\omega_2=HH$ (i.e.\ $\mM_2=2$),
$\mE[\mM_3\mid\text{atom}] = 2+(p-q)$ exactly as before. On the atom
$\omega_1\omega_2=HT$ (i.e.\ $\mM_2=0$), $\mM_3 = 0+\mX_3$ and $\mX_3$
is independent of the atom, so $\mE[\mM_3\mid \text{atom}] = p-q$.
Likewise on the atom $\omega_1\omega_2 = TH$ (also $\mM_2=0$),
$\mE[\mM_3\mid\text{atom}]=p-q$. On the atom $\omega_1\omega_2=TT$
(i.e.\ $\mM_2=-2$), $\mE[\mM_3\mid\text{atom}] = -2+(p-q)$. So
\begin{equation}
  \mE[\mM_3\mid\cA_2](\omega) = \left\{\begin{array}{cl} 2+(p-q) & \omega_1\omega_2 = HH\ ,\\ (p-q) & \omega_1\omega_2 = HT\ ,\\ (p-q) & \omega_1\omega_2 = TH\ ,\\ -2+(p-q) & \omega_1\omega_2 = TT\ . \end{array}\right.
\end{equation}

\underline{Comparison:} Even though $\cA_2$ has strictly finer atoms
than $\sigma(\mM_2)$ (it splits the atom $\{\mM_2=0\}$ into the two
separate atoms $HT$ and $TH$), the value of the conditional
expectation on both of these finer atoms is the same, $p-q$, which
equals the value that $\sigma(\mM_2)$ assigns to the combined atom
$\{\mM_2=0\}$. So, as random variables,
\begin{equation}
  \mE[\mM_3\mid\sigma(\mM_2)] = \mE[\mM_3\mid \cA_2] = \mM_2 + (p-q)\ ,
\end{equation}
they are \emph{identical}, even though they are defined using two
different $\sigma$-algebras, one strictly coarser than the other. The
reason is that the extra information in $\cA_2$ (knowing $\mM_1$, i.e.,
whether the first toss individually was $H$ or $T$, on top of knowing
$\mM_2$) is irrelevant to $\mX_3$, which is independent of the first
two tosses; so this extra information does not change the conditional
expectation of $\mM_3$. \hfill$\blacksquare$

Optional exercises for more practice (will not be graded):

\textbf{Exercise 1i:} Compute all the random variables
$\mE[\mM_m\mid\sigma(\mM_n)]$ for all $1 \leq m,n \leq N=3$. Do you
observe a pattern?


\textbf{Exercise 1j:} Compute all the random variables
$\mE[\mM_m\mid \cA_n]$ for all $1\leq m,n\leq N=3$. How are the
$\sigma$-algebras $\cA_n$ related to the $\sigma$-algebras $\cF_n$
determined by the first $n$ coin tosses that we discussed in class?


\end{document}
